% Both reflections, and the operator they force (paper 14 of the O lane).
% Lane: O (operator bridge), modules YangMills/OS/SpatialOS.lean (site bridge)
% and YangMills/OS/SpatialReconstruction.lean (the operator).
% REGIME: tricotomy; NO scores, NO desk language, NO commit chronology beyond
% the reproducibility freeze.
\documentclass[11pt]{article}
\usepackage[margin=1.1in]{geometry}
\usepackage{amsmath,amssymb,amsthm}
\usepackage{booktabs,longtable,array}
\usepackage[colorlinks=true,linkcolor=blue,urlcolor=blue,citecolor=blue]{hyperref}

\newtheorem{theorem}{Theorem}[section]
\newtheorem{lemma}[theorem]{Lemma}
\newtheorem{proposition}[theorem]{Proposition}
\newtheorem{corollary}[theorem]{Corollary}
\newtheorem{remark}[theorem]{Remark}
\theoremstyle{definition}
\newtheorem{definition}[theorem]{Definition}

\emergencystretch=3em
\tolerance=2000

\newcommand{\anchor}{c90dc745ab8cd1ba7ddc02aa16fed3de339bf958}
\newcommand{\osline}[3]{%
  \href{https://github.com/lluiseriksson/THE-ERIKSSON-PROGRAMME/blob/\anchor/YangMills/OS/#1.lean\#L#2}{\nolinkurl{#3}}}
\newcommand{\repofile}[2]{%
  \href{https://github.com/lluiseriksson/THE-ERIKSSON-PROGRAMME/blob/\anchor/#1}{\nolinkurl{#2}}}
\newcommand{\R}{\mathbb{R}}
\newcommand{\C}{\mathbb{C}}
\newcommand{\Z}{\mathbb{Z}}

\title{Both Reflections, and the Operator They Force\\
\large a machine-checked site reflection bridge, and the self-adjoint positive
transfer operator the two geometries determine, for the coupled $\Z_2$ slice}
\author{Lluis Eriksson}
\date{}

\begin{document}
\maketitle

\begin{abstract}
\textbf{The site bridge, which was the open half.}  A whole path $X$ of $2m+1$
slices through a site plane is exactly a pair of half-paths that \emph{share}
the slice the plane passes through --- a fibred product, not a product --- and
its Gibbs weight satisfies the division-free identity
$w(\sigma)\,W(X)=W(a)\,W(b)$, where $\sigma$ is the shared slice.  That identity
carries no hypothesis on $w$ whatsoever.  With it, the site reflected form of a
complex observable of a whole half-chain is the reflected two-point sum of the
\emph{measure}, at every $\beta$ including negative ones, and its Gram matrix
over a finite family is positive semidefinite.  The companion bond statement was
proved earlier; this closes the pair.

\textbf{The operator the pair forces.}  Both forms factor through the same map
$\Phi$, which sums out the interior of a half and keeps its boundary slice.
$\Phi$ is surjective, so the physical space is the \emph{whole} space of
boundary vectors and nothing smaller.  On it the site form is definite, so if
there is to be an operator $T$ with
$\langle u,Tv\rangle_{\mathrm{site}}=\langle u,v\rangle_{\mathrm{bond}}$ then
$T$ is not chosen but determined: $(Tv)(\sigma)=w(\sigma)\sum_\tau
K(\sigma,\tau)v(\tau)$.  We prove that this $T$ is \textbf{self-adjoint} for the
physical inner product at every $\beta$, and \textbf{positive} for $\beta\ge0$.
Self-adjoint and positive is what an Osterwalder--Schrader transfer operator is
asked to be.

\textbf{Why both geometries had to exist first, and it is not symmetry.}  $T$
advances a half-chain by \emph{one} slice, which turns an even separation into
an odd one.  So the defining equation has the site geometry on its left and the
bond geometry on its right: a construction owning only one of the two has no
equation to define $T$ by at all.

Everything is machine-checked in Lean~4 against Mathlib, with no \texttt{sorry},
no project axioms, and \texttt{propext}, \texttt{Classical.choice},
\texttt{Quot.sound} the only axioms anywhere.  Two independent gate scripts were
committed \emph{before} the corresponding Lean existed and \emph{before} they
were run; each predicts a number rather than a range, and each file states what
its own failure would have meant.

\textbf{What is not here.}  No Hamiltonian, no spectral gap for $T$, no
contraction bound, no infinite volume, no continuum.  The numerical gate
observes that $T$ has the same spectrum as the symmetrised kernel of earlier
work in this lane; that observation is a \emph{gate}, not a theorem, and is not
used anywhere.
\end{abstract}

\section{What is proved, and what is not}\label{sec:scope}

Two things, in the order they were done.

\paragraph{(i) The site bridge.}  The site reflected form of the previous paper
in this lane was a \emph{candidate}: positive, with a positive semidefinite Gram
matrix, but not yet identified with a sum over whole paths against the Gibbs
weight.  Section~\ref{sec:sitebridge} closes it.  The bond case was closed
earlier; the site case is a different statement, because the two halves share a
slice rather than being disjoint.

\paragraph{(ii) The operator.}  Section~\ref{sec:forced} shows the two closed
bridges determine an operator, Section~\ref{sec:sa} proves it self-adjoint and
positive, and Section~\ref{sec:phys} proves the space it acts on is the whole
boundary space rather than a proper subspace of it.

\paragraph{What is \emph{not} proved.}  There is no Hamiltonian here.  There is
no proof that $T$ inherits the Perron eigenvalue or the spectral gap of the
symmetrised kernel studied earlier in this lane; a numerical gate observes that
it does, and an observation registered as a gate is evidence about a plan, not a
theorem, and is used for nothing.  There is no contraction bound.  There is no
thermodynamic limit and no continuum limit.  The system is a finite coupled
$\Z_2$ slice, and every statement below is a statement about it.

\section{The objects}\label{sec:objects}

Fix $L\ge0$ and write a \emph{slice} for a configuration
$\sigma:\{0,\dots,L-1\}\to\{0,1\}$; there are $2^L$ of them.  For $\beta\in\R$
the \emph{slice kernel} is
\[
  K_\beta(\sigma,\tau)\;=\;\prod_{j<L}e^{\beta\,s(\sigma_j,\tau_j)},
  \qquad s(x,y)=+1 \text{ if } x=y,\ -1 \text{ otherwise},
\]
which is symmetric, strictly positive, and defined at every real $\beta$.  A
\emph{source weight} is any $w$ with $w(\sigma)>0$ for all $\sigma$; no further
property of $w$ is assumed anywhere in this paper.

A \emph{path} of $N+1$ slices is $X:\{0,\dots,N\}\to\text{slices}$, and its
Gibbs weight is
\[
  W(X)\;=\;\Big(\prod_{t\le N}w(X_t)\Big)\prod_{t<N}K_\beta(X_t,X_{t+1}),
\]
the unnormalised weight; every statement below is about unnormalised sums, which
is the form the reflection-positivity axiom is about.  Dividing by the partition
function is division by a positive number.

A \emph{past half} is a path $a$ of $m+1$ slices, its \emph{boundary slice} is
$\partial a=a_m$, and an \emph{observable of the half-chain} is any
$F:\{\text{halves}\}\to\C$.  The \emph{collapse} is
\[
  (\Phi F)(\sigma)\;=\;\sum_{\partial a=\sigma}W(a)\,F(a),
\]
which sums out the interior of a half and keeps only the slice at its boundary.
Both reflected forms factor through $\Phi$, and this is the only reason the
operator of Section~\ref{sec:forced} can be written down at all.

\section{The site bridge}\label{sec:sitebridge}

Through a \emph{bond} the reflection plane falls between two slices and the two
halves are disjoint, so pairs of halves are all paths.  Through a \emph{site}
the plane falls on a slice, the two halves \emph{share} it, and the pairs that
come from a path are those agreeing there: a fibred product over the shared
slice.  The weight of that slice is then carried by both halves and appears once
in the path.

\begin{definition}[the assembly]
For halves $a,b$ of $m+1$ slices put
\[
  J(a,b)\;=\;\big(a_0,\dots,a_{m-1},\;b_m,\;b_{m-1},\dots,b_0\big),
\]
a path of $2m+1$ slices.  The blocks have sizes $m$ and $m+1$; that choice, and
not the mirror one, is what makes $2m+1$ a successor so that the weight
functional accepts it.
\end{definition}

\begin{lemma}[the fibred bijection]\label{lem:sitebij}
For each slice $\sigma$, $J$ is a bijection from the pairs of halves with
$\partial a=\partial b=\sigma$ onto the paths of $2m+1$ slices whose middle
slice is $\sigma$.
\end{lemma}

Machine-checked as \osline{SpatialOS}{1051}{sum\_pathsAt\_eq}, with the
two directions
\osline{SpatialOS}{911}{pastSiteOf\_joinSite} and
\osline{SpatialOS}{903}{futSiteOf\_joinSite}.  The first needs the
halves to agree at the boundary and the second does not, which is the
asymmetry the fibred product introduces.

\begin{theorem}[the site weight identity, division-free]\label{thm:siteweight}
If $\partial a=\partial b$ then
\[
  w(\partial a)\;W\big(J(a,b)\big)\;=\;W(a)\,W(b).
\]
No hypothesis on $w$ is used, in particular not positivity.
\end{theorem}

Machine-checked as \osline{SpatialOS}{1017}{gibbsWeight\_joinSite}.  It
splits into a statement about the $w$ factors and one about the kernel factors,
and the two are not equally strong: the $w$ identity
\osline{SpatialOS}{962}{prod\_w\_joinSite} holds for \emph{any} pair of
halves, agreeing at the boundary or not, while the kernel identity
\osline{SpatialOS}{986}{prod\_K\_joinSite} needs them to agree.  We state
this because the first version of the $w$ lemma carried the hypothesis, and the
build's own unused-variable warning was what pointed out that it does not need
it.

\begin{remark}
The kernel identity is where the site case differs from the bond case in
substance rather than in bookkeeping.  Through a bond, exactly one kernel factor
crosses the plane and survives into the weight identity.  Through a site,
\emph{no} kernel factor crosses: the $2m$ bonds of the assembled path are
exactly the $m$ bonds of each half.  The missing crossing factor is the whole
difference, and it is why the site statements hold at every $\beta$ while the
bond ones need $\beta\ge0$.
\end{remark}

\begin{theorem}[the site bridge]\label{thm:sitebridge}
For every $L$, every $m$, every strictly positive $w$ and every $\beta\in\R$,
and for complex observables $F,G$ of a whole half-chain,
\[
  \sum_{X}\ \overline{F(\mathrm{past}\,X)}\ G(\mathrm{rev}\,X)\ W(X)
  \;=\;\sum_\sigma\ \sum_{\partial a=\sigma}\ \sum_{\partial b=\sigma}
    \overline{F(a)}\,G(b)\,\frac{W(a)W(b)}{w(\sigma)} ,
\]
the sum on the left running over all paths of $2m+1$ slices, with
$\mathrm{past}\,X$ its past half and $\mathrm{rev}\,X$ its future half read
backwards from the far end.
\end{theorem}

Machine-checked as
\osline{SpatialOS}{1088}{osPairingSiteCross\_eq\_gibbsSum}; the
diagonal case is \osline{SpatialOS}{1116}{osPairingSite\_eq\_gibbsSum}.
Positivity of $w$ enters in exactly one place --- dividing by $w(\sigma)$ ---
and Theorem~\ref{thm:siteweight} is stated without division so that this is
visible rather than buried.

\begin{corollary}\label{cor:sitepos}
The left-hand side of Theorem~\ref{thm:sitebridge} with $G=F$ is a non-negative
real, for every $\beta\in\R$, and the Gram matrix over a finite family of
observables with complex coefficients is positive semidefinite.
\end{corollary}

\osline{SpatialOS}{1128}{gibbsSumSite\_reflected\_nonneg} and
\osline{SpatialOS}{1141}{gibbsSumSite\_reflected\_gram\_nonneg}.  The
Gram statement is a consequence of the diagonal one applied to
$\sum_i c_iF_i$, and it is written down anyway, because leaving it implicit
would let the prose say ``matrix statement about the measure'' while the formal
text exhibited only a diagonal identity.

\section{The operator is forced}\label{sec:forced}

Write, for boundary vectors $u,v:\{\text{slices}\}\to\C$,
\[
  \langle u,v\rangle_{\mathrm{site}}=\sum_\sigma
    \frac{\overline{u(\sigma)}\,v(\sigma)}{w(\sigma)},
  \qquad
  \langle u,v\rangle_{\mathrm{bond}}=\sum_{\sigma,\tau}
    \overline{u(\sigma)}\,K_\beta(\sigma,\tau)\,v(\tau).
\]
By Theorem~\ref{thm:sitebridge} and its bond companion, evaluated at
$u=\Phi F$, $v=\Phi G$ these are the two reflected forms of the measure ---
\osline{SpatialReconstruction}{91}{siteForm\_collapse} and
\osline{SpatialReconstruction}{98}{bondForm\_collapse}.

The site form is definite as soon as $w>0$.  So the equation
\begin{equation}\label{eq:star}
  \langle u,Tv\rangle_{\mathrm{site}}\;=\;\langle u,v\rangle_{\mathrm{bond}}
  \qquad\text{for all }u
  \tag{$*$}
\end{equation}
has at most one solution $T$, and reading it off gives
\[
  (Tv)(\sigma)\;=\;w(\sigma)\sum_\tau K_\beta(\sigma,\tau)\,v(\tau).
\]

\begin{theorem}[\eqref{eq:star}]\label{thm:star}
With $T$ as displayed, $\langle u,Tv\rangle_{\mathrm{site}}
=\langle u,v\rangle_{\mathrm{bond}}$ for all boundary vectors $u,v$, every
$\beta\in\R$ and every strictly positive $w$.
\end{theorem}

\osline{SpatialReconstruction}{116}{siteForm\_transferOp}.

\paragraph{One geometry on each side.}  $T$ advances a half-chain across one
more slice.  If the pairing on the left is taken at separation $2m$, the pairing
on the right is at separation $2m+1$: the parity flips.  So \eqref{eq:star} is
an equation between a \emph{site} quantity and a \emph{bond} quantity, and it
could not have been written at all with only one of the two bridges in hand.
This is the reason Section~\ref{sec:sitebridge} is in this paper rather than
being a tidy-up of the previous one.

\section{Self-adjointness and positivity}\label{sec:sa}

\begin{theorem}[self-adjoint]\label{thm:sa}
For every $\beta\in\R$, negative included, and every strictly positive $w$,
\[
  \langle u,Tv\rangle_{\mathrm{site}}\;=\;\langle Tu,v\rangle_{\mathrm{site}}.
\]
\end{theorem}

\osline{SpatialReconstruction}{155}{transferOp\_selfAdjoint}, via
\osline{SpatialReconstruction}{130}{siteForm\_transferOp\_left}.  The
left-hand version needs the symmetry of $K_\beta$, which the right-hand version
of \eqref{eq:star} did not; the two are therefore proved separately rather than
one being quoted for the other.

\begin{theorem}[positive]\label{thm:pos}
For $\beta\ge0$ and every strictly positive $w$, $\langle
u,Tu\rangle_{\mathrm{site}}$ is a non-negative real.
\end{theorem}

\osline{SpatialReconstruction}{208}{transferOp\_nonneg}.  By
Theorem~\ref{thm:star} this is the statement that the bond form is positive
semidefinite on complex vectors, which reduces by a generic real-to-complex step
to the positive semidefiniteness of $K_\beta$ proved earlier in this lane.  The
physical inner product itself is positive at every $\beta$:
\osline{SpatialReconstruction}{164}{siteForm\_self\_nonneg}.

\paragraph{Sufficiency only.}  $\beta\ge0$ is proved sufficient for
Theorem~\ref{thm:pos} and is not proved necessary, here or anywhere in this
lane.  At $L=0$ it could not be necessary: the kernel is the scalar $1$.

\section{Definiteness, and why the operator is unique}\label{sec:unique}

Section~\ref{sec:forced} said the defining equation has \emph{at most} one
solution.  That step needs the site form to separate points, and
non-negativity does not separate anything: version 1.0 of this paper asserted
the conclusion while carrying only
\osline{SpatialReconstruction}{SITENNVEC}{siteForm\_self\_nonneg}.  The
missing chain is supplied here.

\begin{theorem}[definiteness]\label{thm:def}
For strictly positive $w$, $\langle u,u\rangle_{\mathrm{site}}=0$ if and only if
$u=0$.
\end{theorem}

\osline{SpatialReconstruction}{DEFLINE}{siteForm\_self\_eq\_zero\_iff}, whence
separation --- two vectors pairing identically against everything are equal ---
at \osline{SpatialReconstruction}{SEPLINE}{siteForm\_right\_ext}.

\begin{theorem}[uniqueness]\label{thm:uniq}
If $S$ is any map with $\langle u,Sv\rangle_{\mathrm{site}}
=\langle u,v\rangle_{\mathrm{bond}}$ for all $u,v$, then $S=T$.
\end{theorem}

\osline{SpatialReconstruction}{UNIQLINE}{transferOp\_unique}.  The statement
quantifies over an arbitrary map, linearity not assumed; linearity of $T$
itself is a separate object,
\osline{SpatialReconstruction}{OPLLINE}{transferOpL}, so that the word
\emph{operator} names a type.

\section{The similarity, and what it does not give}\label{sec:spec}

Let $Q$ be multiplication by $\sqrt{w}$.  For strictly positive $w$ this is a
linear \emph{isomorphism} --- \osline{SpatialReconstruction}{QEQUIVLINE}{sqrtWeightEquiv}
--- and

\begin{theorem}[similarity]\label{thm:sim}
$T\circ Q=Q\circ S$ as an equality of linear maps, where $S$ is the symmetrised
kernel $\sqrt{w}K_\beta\sqrt{w}$ studied earlier in this lane.
\end{theorem}

\osline{SpatialReconstruction}{SIMILARLINE}{transferOpL\_comp\_sqrtWeightEquiv}.

\begin{corollary}\label{cor:eig}
$\lambda$ is an eigenvalue of $S$ if and only if it is an eigenvalue of $T$, the
non-vanishing eigenvector being carried across by $Q$ in both directions.
\end{corollary}

\osline{SpatialReconstruction}{EIGIFFLINE}{transferOp\_eigenvalue\_iff}.

\paragraph{What this is not, stated because the first version of it overstated
itself.}  Gate R4 of \S\ref{sec:gates} asserted four things: equality of
eigenvalue lists entry by entry, reality of all of them, agreement of
multiplicities, and the ratio of the largest modulus to the Perron eigenvalue.
Theorem~\ref{thm:sim} and Corollary~\ref{cor:eig} give the similarity and the
equality of eigenvalue \emph{sets}.  Multiplicities, characteristic
polynomials, ordered lists, reality and the Perron ratio are \textbf{not}
formalised here.  They follow from a similarity with a real symmetric matrix by
standard linear algebra; standard is not the same as present, and an
intermediate version of this paper wrote ``the spectrum transports'' when it had
proved less.

The weaker forward implication is kept and labelled as such at
\osline{SpatialReconstruction}{FWDEIGLINE}{transferOp\_eigen\_of\_symWeighted};
note that it does not require the eigenvector to be non-zero, so at $u=0$ its
hypothesis holds for every $\lambda$ and it says nothing.  That is exactly why
it is not, by itself, a statement about spectra.

\section{The physical space is everything}\label{sec:phys}

An operator statement is only as strong as the space it is about.  If $\Phi$
were not surjective, the physical space would be a proper subspace of the
boundary vectors and every theorem above would be quietly weaker than it sounds.

\begin{theorem}[surjectivity]\label{thm:surj}
For every strictly positive $w$, every $\beta\in\R$ and every $m$, every
boundary vector is $\Phi F$ for some observable $F$ of the half-chain.
\end{theorem}

\osline{SpatialReconstruction}{230}{collapse\_surjective}.  The witness is
explicit: $F(a)=v(\partial a)\big/\sum_{\partial b=\partial a}W(b)$, whose
denominator is positive because the constant half ending at a slice is a half
ending at it --- \osline{SpatialReconstruction}{222}{const\_mem\_halvesAt}
--- and Gibbs weights are positive.

The null space of the site form on observables is exactly $\ker\Phi$ ---
\osline{SpatialReconstruction}{KERLINE}{mem\_ker\_collapseL\_iff}, so the
quotient is by something identified rather than by an abstract null space --- and
the quotient by it \emph{is} the space of boundary vectors, as a bundled linear
equivalence at \osline{SpatialReconstruction}{PHYSLINE}{physicalEquiv}.  Gate R1
of Section~\ref{sec:gates} predicted its dimension as an integer before any of
this was written.  In v1.0 the last two sentences were prose; they are objects
now.

\section{On the measure}\label{sec:measure}

Everything above is about boundary vectors.  Composed with the two bridges it is
about the measure, which is the only reason it is worth stating.

\begin{theorem}\label{thm:onmeasure}
For every strictly positive $w$, every $\beta\in\R$, every $m$ and all complex
observables $F,G$ of a half-chain,
\[
  \big\langle \Phi F,\;T\,\Phi G\big\rangle_{\mathrm{site}}
  \;=\;\sum_{X}\ \overline{F(\mathrm{past}\,X)}\ G(\mathrm{rev}\,X)\ W(X),
\]
the sum running over all paths of $2m+2$ slices.
\end{theorem}

\osline{SpatialReconstruction}{279}{osPairing\_transfer} composed with the
bond bridge, and stated on the path sum itself at
\osline{SpatialReconstruction}{294}{osPairing\_transfer\_gibbsSum}.  In
words: the matrix element of the transfer operator between two half-chain
observables \emph{is} the reflected two-point function of the Gibbs measure one
slice further apart.

\begin{remark}[why the cross form exists]
The displayed identity is stated for two observables, and the formal text states
it for two.  It would have been cheaper to prove only the diagonal case $G=F$
and let the display carry the general one; the bond bridge's cross form had in
fact been proved inside another proof, where nothing could cite it, and was
extracted --- \osline{SpatialOS}{785}{osPairingBondCross\_eq\_gibbsSum}
--- precisely so that this section's prose and its formal text have the same
scope.  A green build does not notice a theorem carrying less than its sentence.
\end{remark}

\section{The gates}\label{sec:gates}

Two scripts, each committed before the Lean it licenses existed and before it
was run, each predicting a \emph{number} rather than a range, and each stating
in its own text what its failure would have meant.

\paragraph{Site bridge, \repofile{scripts/judge\_site\_bridge.py}{scripts/judge\_site\_bridge.py}.}
\textbf{S1} predicted that the count of paths and the count of assembling pairs
are the same integer, $2^{L(2m+1)}$, in every cell.  \textbf{S2} predicted the
residual of Theorem~\ref{thm:siteweight} to be zero to $10^{-12}$ relative.  The
file states that an S2 residual which was a fraction of the weight rather than
roundoff would have meant the site form was the \emph{wrong} form, not merely an
unidentified one.  Both passed.

\paragraph{Reconstruction, \repofile{scripts/judge\_os\_reconstruction.py}{scripts/judge\_os\_reconstruction.py}.}
\textbf{R1} predicted $\operatorname{rank}\Phi=2^L$ exactly, as an integer.
\textbf{R2} predicted the self-adjointness residual to be zero to $10^{-12}$,
negative $\beta$ included.  \textbf{R3} predicted the residual of \eqref{eq:star}
itself to be zero.  \textbf{R4} predicted that the eigenvalues of $T$ agree
entry by entry with those of the symmetrised kernel $\sqrt{w}K_\beta\sqrt{w}$,
that all are real, and that the largest in modulus divided by the Perron
eigenvalue is exactly $1$.  All four passed.

\textbf{R4 is a gate and is used for nothing.}  It is recorded because it says
where the next theorem would come from, and because registering it before the
fact is the only way it can later count as evidence at all.  It is not cited in
any proof and no statement in this paper depends on it.

Neither script lets an acceptance decision depend on a Python \texttt{assert};
both are run in the ordinary and the \texttt{-O} interpreter modes, because
\texttt{-O} deletes assertions and this repository has twice emitted a false
\textsc{pass} that way.

\section{What is open}\label{sec:open}

\begin{enumerate}
\item \textbf{The spectrum of $T$.}  R4 observes that $T$ is similar to the
symmetrised kernel, which earlier work in this lane analyses --- Perron vector,
strict spectral gap at every finite extent, an exact ratio in a decoupled case.
Nothing of that is transported here.  Until it is, $T$ is self-adjoint and
positive and nothing is known about its norm.
\item \textbf{No contraction bound.}  Consequently there is no $\|T\|\le\lambda$
statement, and therefore no $H=-\log(T/\lambda)$ and no Hamiltonian.
\item \textbf{No thermodynamic limit.}  Every constant here may depend on $L$
and on $m$.  Uniformity in the extent is a separate and open question for the
coupled kernel.  We note only that with a ring source weight the symmetrised
kernel \emph{is}, by comparing the two definitions, the transfer matrix of the
anisotropic square-lattice Ising model --- an identification, not a cited
theorem --- which is a reason to expect the question to be hard rather than a
result about it.
\item \textbf{Necessity of $\beta\ge0$} for Theorem~\ref{thm:pos} is not proved.
\item \textbf{$\Z_2$ only.}  Nothing here is a gauge theory.
\end{enumerate}

\section{Erratum to v1.0}\label{sec:erratum}

Version 1.0 of this paper was published with a wrong number in it, and the
instrument that produced the number is the reason nobody noticed.  Both are
recorded here rather than silently corrected, because a counter that changes
between versions without an explanation is exactly as untrustworthy as one that
does not change when it should.

\paragraph{What was wrong.}  The table in \S\ref{sec:repro} of v1.0 printed
\emph{Declarations in} \texttt{SpatialOS.lean} $=65$.  The true number at that
same source anchor is $69$.

\paragraph{Why.}  The counter was a line-oriented recogniser that did not accept
an attribute block before the declaration keyword, so every
\verb|@[simp] theorem| was invisible to it.  Four declarations were of that
form.  Recomputing the v1.0 blobs with a corrected recogniser gives $69$ and
$15$, against the printed $65$ and $15$; the second module contained no
attributed declaration, which is why only one of the two numbers moved and why
the error was not self-evident.

\paragraph{The consequence that was worse than the number.}  The same blindness
was what checked that every declaration appears in the oracle.  Two of the four
invisible declarations, \verb|pastOf_joinBond| and \verb|futRevOf_joinBond|,
were \emph{genuinely absent} from the oracle --- inherited from the previous
paper of this lane, and never noticed because the instrument that would have
noticed could not see them.  They are present now, and the count of declarations
carrying an axiom report is reconciled for \emph{both} modules independently
rather than for the new one only.

\paragraph{What was not wrong.}  The build was green and the axiom report was
clean at v1.0; those rows of its table stand.  The defect was in a counter and
in the reach of a check, not in an elaborated statement.

\paragraph{What changed so that it cannot recur in this shape.}  The recogniser
now lives in one file rather than in two that had already drifted apart, it
carries a battery of syntactic cases that runs standalone and fails on
regression, and the coverage check keeps the two modules as independent
quantities.  The battery earned itself immediately: the first corrected
recogniser read the sentence \emph{``The Osterwalder--Schrader axiom asks for
more''} out of this paper's own module header and reported a declaration named
\texttt{asks}.  A parser that has once certified a false coverage should not
remain an unexamined root of trust, and the durable fix --- enumerating
declarations from the elaborated Lean environment instead of from source text
--- is \emph{not} done here.

\section{Reproducibility}\label{sec:repro}

All \emph{code and verification} artifacts --- the two modules, the oracle, the
two gate scripts, the ledger --- are at source checkpoint \texttt{\anchor}.
This manuscript and its PDF are committed on top of that checkpoint, since a
document cannot cite the commit that introduces it.  Every hyperlink was
resolved against \texttt{\anchor} and checked to point at an actual declaration
line, before compilation, again inside the compiled PDF, and again live over
HTTP.

Everything was elaborated on \emph{one} fresh Linux clone at that checkpoint,
with the toolchain and the Mathlib revision pinned in the repository and
re-checked before any compilation started.  A second independent reproduction is
\emph{not} claimed: one clone is one clone.  The two module files and the oracle
file were carried back to the desktop and their SHA-256 digests compared against
the ones computed in the build environment; all three match, so the text linked
below is the text that was elaborated.

\begin{center}
\begin{tabular}{ll}
\toprule
Quantity & Value \\
\midrule
Full core build, this version & JOBSAFTER jobs, success \\ % @@JOBSAFTER@@
\quad same tree at the v1.0 anchor & JOBSBEFORE jobs \\ % @@JOBSBEFORE@@
\quad delta, registered before measuring & JOBSDELTA \\ % @@JOBSDELTA@@
\quad at this campaign's parent commit & JOBSCAMPAIGNBASE jobs \\ % @@JOBSCAMPAIGNBASE@@
Oracle commands, and answers & ORACLETOTAL \\ % @@ORACLETOTAL@@
\quad with a non-standard axiom & ORACLENONSTD \\ % @@ORACLENONSTD@@
Declarations in \texttt{SpatialOS.lean} & DECLSOS \\ % @@DECLSOS@@
\quad of them in the oracle & DECLSOSORACLE \\ % @@DECLSOSORACLE@@
Declarations in \texttt{SpatialReconstruction.lean} & DECLSREC \\ % @@DECLSREC@@
\quad of them in the oracle & DECLSRECORACLE \\ % @@DECLSRECORACLE@@
Distinct axioms across all outputs & \texttt{propext}, \texttt{Quot.sound},
  \texttt{Classical.choice} \\
Occurrences of \texttt{sorryAx} & SORRYCOUNT \\ % @@SORRYCOUNT@@
Project axioms & 0 \\
Errors & COREERRORS \\ % @@COREERRORS@@
\bottomrule
\end{tabular}
\end{center}

The delta of exactly $+1$ is this repository's standing check that a new module
is really being elaborated: one module, one job.  The prediction was registered
before any count was taken, and as a \emph{delta} rather than an absolute,
because a count copied from a document is not a baseline.  Both counts in the
table come from core builds in the \emph{same} build environment.

\paragraph{Which base, and a number that first looked like a failure.}  The
first comparison was run against the merge-base with the repository's main
branch and returned $8465\to8469$, a delta of $+4$.  That is not this campaign:
the branch it was measured from also carries an unrelated line of work that had
already added three modules to the core.  Measured against the parent of this
campaign's first commit --- which is what the prediction was about --- the core
is $8468$ and the delta is $+1$.  Both numbers are true and they answer
different questions.  We record the episode rather than only the number that
agreed, because the correct first reaction to a disagreeing measurement is to
suspect it of measuring the wrong thing, not to adjust the claim.

The modules are
\repofile{YangMills/OS/SpatialOS.lean}{YangMills/OS/SpatialOS.lean} and
\repofile{YangMills/OS/SpatialReconstruction.lean}{YangMills/OS/SpatialReconstruction.lean}.

\small
\paragraph{Previous papers of this lane cited above.}  The bond bridge, the site
reflected form as a candidate, and the collapse are from \emph{The Collapse}
(paper 13 of this lane); the reflection-positivity of the endpoints and the
positive semidefiniteness of $K_\beta$ are from the reflection paper (12); the
Perron vector, the strict spectral gap and the exact ratio referred to in
\S\ref{sec:open} are from papers 7, 8 and 11.  None of those results is
re-proved here and none of them is assumed beyond what is cited by name in the
formal text.

\end{document}
